% =========================================================
% Ordered Pairs and Relations
% =========================================================

\section{Ordered Pairs and Relations}

% ---------------------------------------------------------
% TOOLKIT
% ---------------------------------------------------------

\begin{definitionbox}{Definition (Ordered Pair)}
\begin{definition}[Ordered Pair]\label{def:ordered-pair}
Let $a$ and $b$ be sets. The \emph{ordered pair} $(a,b)$ is defined (Kuratowski)
as
\[
(a,b) \;:=\; \{\{a\}, \{a,b\}\}.
\]
\end{definition}
\end{definitionbox}

\begin{theorembox}{Theorem (Uniqueness of Ordered Pairs)}
\begin{theorem}[Uniqueness of Ordered Pairs]\label{thm:ordered-pair-unique}
For any sets $a,b,c,d$,
\[
(a,b) = (c,d)
\;\;\Longleftrightarrow\;\;
(a = c \land b = d).
\]
\hyperref[prf:ordered-pair-unique]{Go to proof.}
\end{theorem}
\end{theorembox}

\NoLocalDependencies

\begin{remark*}[Standard quantified statement]
$\forall a, b, c, d:\; (a,b) = (c,d) \;\leftrightarrow\; (a = c \wedge b = d)$.
\end{remark*}
\begin{remark*}[Kuratowski encoding]
The purpose of the Kuratowski definition is purely to guarantee the uniqueness
theorem above. Once $(a,b) = (c,d) \iff a=c \land b=d$ is established,
we may treat ordered pairs as a primitive with this property and forget the
encoding.
\end{remark*}

\begin{definitionbox}{Definition (Cartesian Product — as foundation for relations)}
\begin{definition}[Cartesian Product — as foundation for relations]\label{def:cartesian-product-rel}
The \emph{Cartesian product} of sets $A$ and $B$ is:
\[
A \times B
\;:=\;
\{\, (a,b) \mid a \in A \text{ and } b \in B \,\}.
\]
\end{definition}
\end{definitionbox}

\begin{definitionbox}{Definition (Relation)}
\begin{definition}[Relation]\label{def:relation}
A \emph{relation} from $A$ to $B$ is any subset $R \subseteq A \times B$.

If $(a,b) \in R$ we write $a \, R \, b$ and say ``$a$ is related to $b$.''
When $A = B$ we say $R$ is a \emph{relation on $A$}.
\end{definition}
\end{definitionbox}

\begin{remark*}[Relations as structured sets]
Relations introduce no new foundational objects: they are simply sets of
ordered pairs, constructed using operations already available. Properties of
relations can therefore be studied with ordinary set-theoretic reasoning.
\end{remark*}
